\SourceLeaf{243}%
bis ich mich in Ruhe wieder diesem Fach widmen kann und mir das Lesen russischer Bücher leichter fällt als jetzt. Wenn ich mich nicht irre, steht auch ein Aufsatz des p. \textsc{Lobatschewski}, vielleicht eine Übersetzung aus den \gausscyr{Записки}, in \textsc{Crelle}s Journal. Das kann ich aber im Augenblick nicht nachsehen. \mbox{.\ .\ .\ .\ .\ .\ .}

\par\bigskip\Needspace{6\baselineskip}

\begin{center}\textsc{Gauss} an \textsc{Gerling}. Göttingen, 8. Februar 1844.\end{center}

Mein liebster Freund, vielleicht ist es Ihnen nicht unlieb, wenn ich die bibliografischen Angaben aus meinem letzten Brief noch um den einen oder anderen Punkt ergänze. Auf die Sache selbst kann ich derzeit freilich nicht näher eingehen.

Der Aufsatz von \textsc{Lobatschewski} in \textsc{Crelle}s Journal steht in Band 17, S. 295 ff. Ich stelle fest, dass er nur eine freie Übersetzung des russischen Aufsatzes [\gausscyr{Воображаемая геометрія}] im Jahrgang 1835 der Gelehrten Schriften der Kasaner Universität ist. Dort stößt man auf dieselben Schwierigkeiten wie im deutschen Aufsatz. In diesem stoßen Sie auf S. 296, Zeile 10, auf die Worte \textls{J’ai démontré} \emph{etc.}, mit denen einem Leser, der nur diesen Aufsatz hat, wenig gedient ist. Dasselbe gilt für \textls{J’ai prouvé ailleurs} \emph{etc.} oben auf S. 303. Der frühere Aufsatz, auf den sich dies zu beziehen \textls{scheint}, dürfte derjenige sein, der in einer Anmerkung des erwähnten russischen Aufsatzes unter dem Titel Über die Anfänge oder Principe der Geometrie im \gausscyr{Казанскомъ Вѣстникѣ} (Kasanschen Boten) für 1829 und 1830 angeführt wird. Dabei wird auch bemerkt, dass eine sehr verletzende Kritik dieser Abhandlung in Nr. 41 der anderen russischen Zeitschrift »Sohn des Vaterlandes« \gausscyr{Сынъ Отечества} von 1834 steht, die, wie ich vermute, in Petersburg erscheint. \textsc{Lobatschewski} hatte dagegen eine Erwiderung eingesandt, die aber bis Anfang 1835 nicht aufgenommen worden war.

Mit diesen bibliografischen Angaben ist uns freilich wenig geholfen, da sich in Deutschland kaum ein Exemplar des Kasaner Boten von\space%
